% Word count aim: 750-1000
\chapter{Conclusion and Future Work}
\label{cha:conc}

\section{Conclusion}

This project set out to address a gap that \citet{pinto2018practice}
identified explicitly but left open: the absence of a \textit{combined
anytime measure} capable of describing the full quality trajectory of a
black-box optimization algorithm rather than only its terminal behavior.
The three success criteria established in
Section~\ref{subsec:intro:success} provided the evidentiary floor
against which that gap could be declared closed.

\paragraph{Theoretical criterion.}
Theorem~\ref{thm:regret-profile} establishes the required link between
expected incumbent \acrfull{cr} and the distributional runtime profile
$\mathcal{P}(v,t) = \mathbb{P}(\tau_v > t)$. Together with
Corollary~\ref{cor:asymptotic-regret}, it shows that long-run
\acrshort{cr} is governed by expected first-hitting times across
all fitness levels, not only by the time to the optimum. This supplies
the unified anytime measure called for in the literature.

The criterion further required numerical corroboration within
Monte-Carlo sampling error on at least two benchmark functions. Suite
08 (Section~\ref{sec:results:suite08}) provides this corroboration on
\textsc{OneMax} and \textsc{LeadingOnes} for both \acrshort{rls} and
the $(1+1)$-\acrshort{ea}, with scale-normalized maximum pointwise
errors (\acrshort{rmae}) below $1\%$ in all four combinations
(Table~\ref{tab:profile-verification-distances}).
The theoretical criterion is therefore met for the integer-valued benchmark
setting tested here.

\paragraph{Empirical criterion.}
The empirical criterion required \acrshort{cr} trajectories to
yield qualitatively distinct signatures across at least three landscape
archetypes that classical measures cannot separate. The experimental
evidence meets this standard on multiple axes.

On smooth unimodal landscapes (\textsc{OneMax}, \textsc{LeadingOnes},
\textsc{BinVal}), Suites~01 and~02 show that incumbent \acrshort{cr}
curves converge to finite asymptotes whose ordering can \textit{invert}
the \acrshort{ttfo} ordering. Thus, the asymptotic identity has direct
empirical consequence: early progress across intermediate fitness
levels matters, even when the final optimum is eventually reached.

On barrier and deceptive landscapes (\textsc{Jump}, \textsc{Trap},
\textsc{Plateau}), where $P(\text{opt}) \approx 0$ within the
available budget and \acrshort{ttfo} distributions provide no usable
samples, \acrshort{cr} trajectories remain informative. Even when
barrier archetypes are not perfectly separated within the finite
budget, their slopes expose residual fitness gaps and stagnation
levels. This is the operational content of the project's central claim:
regret provides structured information precisely where binary
convergence measures go silent.

SA cooling schedule aggressiveness is directly legible in the
instantaneous regret trajectory (Suite~11), with the onset of the
low-regret floor ordering schedules from fastest to slowest. On NK
landscapes, the regret plots expose the cooling--ruggedness trade-off:
temperature-driven exploration is costly on low-epistasis instances but
becomes competitive, and in calibrated cases beneficial, as epistasis
increases. The empirical criterion is met across more than three
landscape archetypes.

\paragraph{Software criterion.}
The \texttt{regret} framework provides the required reproducible
execution layer for all reported suites, including trajectory
recording, regret computation, runtime-profile analysis, and
\acrfull{ci}-enforced branch coverage across Python 3.12, 3.13, and
3.14.

All three success criteria are therefore satisfied. In answering
\citet{pinto2018practice}'s call, this project demonstrates that
incumbent \acrshort{cr} is a natural extension of fixed-budget
analysis. A connection was established between \acrshort{cr}
and runtime profiles on integer-valued pseudo-Boolean landscapes, with
discretization required beyond that setting.

\section{Future Work}

\paragraph{Tight analytical and asymptotic regret bounds.}
While Theorem~\ref{thm:regret-profile} provides a route from known
hitting-time bounds to regret bounds via the distributional runtime
profile, the resulting expressions are in general loose. Deriving
tight asymptotic characterizations of $\mathbb{E}[\mathrm{CR}(T)]$ for
\acrshort{rls} and the $(1+1)$-\acrshort{ea} on
\textsc{OneMax}, \textsc{LeadingOnes}, and $\textsc{Jump}_k$ remains an
open theoretical problem. A larger scaling study over increasing $n$
would complement this by testing whether observed regret asymptotes
grow at the rates predicted by summing known $\mathbb{E}[\tau_v]$
expressions across fitness levels.

\paragraph{Extension to unknown optima.}
A key assumption throughout this project is that the global optimum
$f^*$ is known in advance, which is necessary for computing exact
regret values. This assumption holds for the emulated benchmark
problems studied here but rarely in genuine black-box practice. Future
work could pursue two complementary directions: the development of
principled heuristics for estimating $f^*$ online, and the adoption of
a probabilistic treatment of the global optimum that would allow the
tail-sum identity to be applied under uncertainty. Problems where the
optimum is known apriori, but the objective landscape remains hidden
represent a particularly tractable intermediate case: the regret
framework developed here provides the necessary foundation for
exploiting that knowledge.

\paragraph{Regret-optimal parameter schedules.}
The $\varepsilon$-greedy and decaying-$\varepsilon$-greedy
\acrshort{rls} variants studied in this project use exploration
schedules motivated by classical bandit theory, but it is not known
whether these schedules minimize \acrshort{cr} in the
pseudo-Boolean setting, where reward structure is non-stationary.
An analog of the UCB1 construction~\cite{auer2002finite} adapted to
the search-space structure of $\{0,1\}^n$ could provide a principled
route to parameter control with explicit regret guarantees, connecting
the exploration--exploitation analysis of online learning to the
adaptive operator selection literature in evolutionary computation.

\paragraph{High-probability tail bounds.}
The analyses presented here concern expected quantities, corresponding
to the Las Vegas framework. Deriving high-probability tail bounds on
incumbent \acrshort{cr}, analogous to the concentration results
available for bandit algorithms~\cite{lattimore2020bandit}, would
allow practitioners to give finite-confidence guarantees on worst-case
trajectory quality rather than only on its expectation.
